\section{广义 Verma 模的 Kazhdan--Lusztig 算法}\label{sec:06.6.3}

设 $\Theta\subset\Pi$。本节假设 $\lambda=-\rho$，并记
$$I_O=I(O,-\rho),\qquad L_O=L(O,-\rho).$$

考虑 V.2 中引入的 $\Z[q,q^{-1}]$-模 $H$。对 $W$ 的每个右 $W_\Theta$-陪集 $C$，记 $w_C$ 为该陪集中的最长元，并记 $\delta_C$ 为 $H$ 中由
$$\delta_C=\sum_{v\in W_\Theta}q^{\ell(v)}\delta_{vw_C}$$
给出的元素。设 $H_\Theta$ 是 $H$ 中由诸 $\delta_C$（$C\in W_\Theta\backslash W$）张成的 $\Z[q,q^{-1}]$-子模。

设 $C\in W_\Theta\backslash W$、$\alpha\in\Pi$。由 1.6，只有三种可能：$Cs_\alpha=C$；$Cs_\alpha>C$；$Cs_\alpha<C$。下面在这些情形下计算 $T_\alpha$ 作用在 $\delta_C$ 上的结果。

\begin{lemma}\label{lem:06.3.1}
设 $C\in W_\Theta\backslash W$、$\alpha\in\Pi$。则：
\begin{enumerate}[label=(\roman*)]
\item 若 $Cs_\alpha=C$，则 $T_\alpha(\delta_C)=(q+q^{-1})\delta_C$；
\item 若 $Cs_\alpha>C$，则 $T_\alpha(\delta_C)=q\delta_C+\delta_{Cs_\alpha}$；
\item 若 $Cs_\alpha<C$，则 $T_\alpha(\delta_C)=q^{-1}\delta_C+\delta_{Cs_\alpha}$。
\end{enumerate}
\end{lemma}

\begin{proof}
先看情形 (i)。此时左乘 $s_\alpha$ 置换 $C$ 中的元。置
$$C_+=\{w\in C\mid\ell(ws_\alpha)=\ell(w)+1\},\qquad
C_-=\{w\in C\mid\ell(ws_\alpha)=\ell(w)-1\}.$$
则 $C_+s_\alpha=C_-$，$C_-s_\alpha=C_+$。于是
{\small
\begin{align*}
T_\alpha(\delta_C)&=\sum_{v\in W_\Theta}q^{\ell(v)}T_\alpha(\delta_{vw_C})\\
&=\sum_{vw_C\in C_+}q^{\ell(v)}\bigl(q\delta_{vw_C}+\delta_{vw_Cs_\alpha}\bigr)
+\sum_{vw_C\in C_-}q^{\ell(v)}\bigl(q^{-1}\delta_{vw_C}+\delta_{vw_Cs_\alpha}\bigr).
\end{align*}
}
如前所述，若 $vw_C\in C_+$，则 $vw_Cs_\alpha\in C_-$，故 $vw_Cs_\alpha=v'w_C$，其中
$$\ell(v')=\ell(w_C)-\ell(v'w_C)=\ell(w_C)-\ell(vw_Cs_\alpha)=\ell(w_C)-\ell(vw_C)-1=\ell(v)-1.$$
同理，若 $vw_C\in C_-$，则 $vw_Cs_\alpha\in C_+$，且 $vw_Cs_\alpha=v'w_C$，其中 $\ell(v')=\ell(v)+1$。因此
{\small
\begin{align*}
T_\alpha(\delta_C)&=\sum_{vw_C\in C_+}q^{\ell(v)+1}\delta_{vw_C}
+\sum_{v'w_C\in C_-}q^{\ell(v')+1}\delta_{v'w_C}\\
&\quad+\sum_{vw_C\in C_-}q^{\ell(v)-1}\delta_{vw_C}
+\sum_{v'w_C\in C_+}q^{\ell(v')-1}\delta_{v'w_C}\\
&=\sum_{v\in W_\Theta}(q+q^{-1})q^{\ell(v)}\delta_{vw_C}=(q+q^{-1})\delta_C.
\end{align*}
}
在情形 (ii)，由 1.6 有 $w_{Cs_\alpha}=w_Cs_\alpha$。因此
{\small
$$T_\alpha(\delta_C)=\sum_{v\in W_\Theta}q^{\ell(v)}T_\alpha(\delta_{vw_C})
=\sum_{v\in W_\Theta}q^{\ell(v)}\bigl(q\delta_{vw_C}+\delta_{vw_Cs_\alpha}\bigr)
=q\delta_C+\delta_{Cs_\alpha}.$$}
在情形 (iii)，同样由 1.6 有 $w_{Cs_\alpha}=w_Cs_\alpha$。因此
{\small
$$T_\alpha(\delta_C)=\sum_{v\in W_\Theta}q^{\ell(v)}\bigl(q^{-1}\delta_{vw_C}+\delta_{vw_Cs_\alpha}\bigr)
=q^{-1}\delta_C+\delta_{Cs_\alpha}.$$}
\end{proof}

\begin{corollary}\label{cor:06.3.2}
$H_\Theta$ 在诸 $T_\alpha$（$\alpha\in\Pi$）下不变。
\end{corollary}

\begin{lemma}\label{lem:06.3.3}
设 $C\in W_\Theta\backslash W$。则 $\phi(w_C)\in H_\Theta$，且
$$\phi(w_C)=\delta_C+\sum_{D<C}P_{w_Cw_D}\delta_D.$$
\end{lemma}

\begin{proof}
设 $O$ 是对应于 $w_C$ 的 $P_\Theta$-轨道。则由 2.3 知 $L_{w_C}=L_O$。设 $O'$ 是 $\overline O$ 中的 $P_\Theta$-轨道，$i_{O'}$ 是 $O'$ 到 $X$ 的自然浸入。则 $R^pi_{O'}^!(L_O)$ 是 $O'$ 上的 $P_\Theta$-等变联络，即若干 $\calO_{O'}$ 的直和。设 $D\in W_\Theta\backslash W$ 对应于 $O'$。则
$$O'=\bigcup_{v\in W_\Theta}C(vw_D).$$
又对任意 $v\in W_\Theta$，$C(vw_D)$ 是 $O'$ 的光滑子簇，余维为 $\ell(v)$。因此，若记 $j$ 为 $C(vw_D)$ 到 $O'$ 的浸入，则当 $p\neq\ell(v)$ 时 $R^pj_!(\calO_{O'})=0$，且
$$R^{\ell(v)}j_!(\calO_{O'})=\calO_{C(vw_D)}.$$
这说明，对任意 $p\in\Z$ 有
$$R^pi_{vw_D}^!(L_O)=R^{\ell(v)}j_!\bigl(R^{p-\ell(v)}i_{O'}^!(L_O)\bigr).$$
因此
$$\dim R^pi_{w_D}^!(L_O)=\dim R^pi_{O'}^!(L_O),$$
且对任意 $p\in\Z$、$v\in W_\Theta$ 有
$$\dim R^pi_{vw_D}^!(L_O)=\dim R^{p-\ell(v)}i_{O'}^!(L_O)=\dim R^{p-\ell(v)}i_{w_D}^!(L_O).$$
这说明
{\small
\begin{align*}
\phi(w_C)&=\nu(L_{w_C})=\nu(L_O)=\sum_{w\in W}\sum_{m\in\Z}\dim\bigl(R^mi_w^!(L_O)\bigr)q^m\delta_w\\
&=\sum_{v\in W_\Theta}\sum_{D\in W_\Theta\backslash W}\sum_{m\in\Z}\dim\bigl(R^{m-\ell(v)}i_{w_D}^!(L_O)\bigr)q^m\delta_{vw_D}\\
&=\sum_{v\in W_\Theta}\sum_{D\in W_\Theta\backslash W}\sum_{p\in\Z}\dim\bigl(R^pi_{w_D}^!(L_O)\bigr)q^{p+\ell(v)}\delta_{vw_D}\\
&=\sum_{D\in W_\Theta\backslash W}\sum_{p\in\Z}\dim\bigl(R^pi_{w_D}^!(L_O)\bigr)q^p\delta_D
=\sum_{D\in W_\Theta\backslash W}P_{w_Cw_D}\delta_D.
\end{align*}
}
\end{proof}

以下对 $C\in W_\Theta\backslash W$ 置
$$\phi(C)=\phi(w_C).$$

最后我们指出下述事实。

\begin{lemma}\label{lem:06.3.4}
设 $C\in W_\Theta\backslash W$、$\alpha\in\Pi$ 满足 $Cs_\alpha<C$。则
$$T_\alpha(\phi(Cs_\alpha))=\sum_{D\le C}c_D\phi(D),$$
其中 $c_D\in\Z$。
\end{lemma}

\begin{proof}
由 3.2 与 3.3 知 $T_\alpha(\phi(Cs_\alpha))\in H_\Theta$，即
{\small
$$T_\alpha(\phi(Cs_\alpha))=\sum_{D\in W_\Theta\backslash W}Q_D\delta_D
=\sum_{v\in W_\Theta}\sum_{D\in W_\Theta\backslash W}Q_Dq^{\ell(v)}\delta_{vw_D},$$}
其中 $Q_D\in\Z[q,q^{-1}]$。另一方面，由 1.6 知 $w_{Cs_\alpha}=w_Cs_\alpha$，且 $\ell(w_{Cs_\alpha})=\ell(w_C)-1$，故由 V.2.1 有
$$T_\alpha(\phi(Cs_\alpha))=T_\alpha(\phi(w_Cs_\alpha))=\sum_{v\le w_C}c_v\phi(v).$$
这说明诸 $Q_D$ 都在 $\Z[q]$ 中。把这些表达式在 $0$ 处取值，得到
$$\sum_{D\in W_\Theta\backslash W}Q_D(0)\delta_{w_D}=T_\alpha(\phi(Cs_\alpha))(0)=\sum_{v\le w_C}c_v\delta_v.$$
这说明 $c_v\neq0$ 蕴含 $v=w_D$，其中 $D\in W_\Theta\backslash W$。于是 $\phi(v)=\phi(D)$，且 $w_D\le w_C$，即 $D\le C$。
\end{proof}

这给出 V.2.1 的下述推广。

\begin{theorem}\label{thm:06.3.5}
存在唯一的函数 $\phi:W_\Theta\backslash W\to H_\Theta$，满足下述性质：
\begin{enumerate}[label=(\roman*)]
\item 对 $C\in W_\Theta\backslash W$ 有
$$\phi(C)=\delta_C+\sum_{D<C}P_{CD}\delta_D,$$
其中 $P_{CD}\in q\Z[q]$；
\item 对 $\alpha\in\Pi$ 与满足 $Cs_\alpha\neq C$、$\ell(w_{Cs_\alpha})=\ell(w_C)-1$ 的 $C\in W_\Theta\backslash W$，存在依赖于 $\alpha$ 与 $C$ 的 $c_D\in\Z$，使
$$T_\alpha(\phi(Cs_\alpha))=\sum_{D\le C}c_D\phi(D).$$
\end{enumerate}
这些多项式 $P_{CD}$ 由 $(W,S)$ 的 Kazhdan--Lusztig 多项式给出：
$$P_{CD}=P_{w_Cw_D}\qquad(C,D\in W_\Theta\backslash W,\ D\le C).$$
\end{theorem}

\begin{proof}
存在性已经证明。剩下证唯一性。这一段论证与 V.2.1 中唯一性的证明类似。

首先可设 $\Theta\neq\Pi$，因为 $\Theta=\Pi$ 时 $W_\Theta=W$，证明是平凡的。

对 $\ell(w_C)$ 作归纳。函数 $C\mapsto\ell(w_C)$ 在 $C=W_\Theta$ 处取到最小值。显然，(i) 给出 $\phi(W_\Theta)=\delta_{W_\Theta}$，而 (ii) 在此情形下是空条件。

取 $C\in W_\Theta\backslash W$ 使 $\ell(w_C)>\ell(w_\Theta)$。由归纳假设，$\phi$ 在一切满足 $\ell(w_D)<\ell(w_C)$ 的 $D\in W_\Theta\backslash W$ 上已唯一确定。于是由 1.7，可取单根 $\alpha$ 使 $Cs_\alpha<C$。由 1.6，此时 $\ell(w_{Cs_\alpha})=\ell(w_C)-1$。由 (ii) 知
$$T_\alpha(\phi(Cs_\alpha))=\sum_{D\le C}c_D\phi(D),$$
且在 $q=0$ 处取值并用 (i) 得
$$T_\alpha(\phi(Cs_\alpha))(0)=\sum_{D\le C}c_D\delta_D.$$
由归纳假设，左端已唯一确定。这说明诸 $c_D$ 唯一确定。另一方面，置 $C'=Cs_\alpha$，有
{\small
\begin{align*}
T_\alpha(\phi(Cs_\alpha))&=T_\alpha\Bigl(\delta_{C'}+\sum_{D<C'}P_{C'D}\delta_D\Bigr)
=T_\alpha(\delta_{C'})+\sum_{D<C'}P_{C'D}T_\alpha(\delta_D)\\
&=q\delta_{C'}+\delta_C+\sum_{D<C'}P_{C'D}T_\alpha(\delta_D).
\end{align*}
}
按构造，$\ell(w_D)<\ell(w_{C'})=\ell(w_C)-1$。故 $T_\alpha(\delta_D)$ 的展开式只能涉及满足 $\ell(w_{D'})\le\ell(w_C)-1$ 的 $\delta_{D'}$，特别地不能涉及 $\delta_C$。这说明 $c_C=1$。于是
$$\phi(C)=T_\alpha(\phi(Cs_\alpha))-\sum_{D<C}c_D\phi(D),$$
这就证明了 $\phi(C)$ 的唯一性。
\end{proof}

\begin{theorem}\label{thm:06.3.6}
设 $P_{CD}$（$C,D\in W_\Theta\backslash W$）是 3.5 中的那些多项式。对 $C\in W_\Theta\backslash W$，记 $O_C$ 为相应的 $P_\Theta$-轨道。则
$$\chi(L_{O_C})=\chi(I_{O_C})+\sum_{D<C}P_{CD}(-1)\chi(I_{O_D}).$$
\end{theorem}

\begin{proof}
因 $I_{O_C}$ 以 $L_{O_C}$ 为唯一的不可约子模，而其余合成因子都是 $L_{O_D}$（$D<C$），可见诸 $\chi(I_{O_C})$（$C\in W_\Theta\backslash W$）构成 $K(\mathcal{M}_{coh}(\D_X,P_\Theta))$ 的一组基。故
$$\chi(L_{O_C})=\sum_{D\le C}\lambda_{CD}\chi(I_{O_D}),$$
其中 $\lambda_{CD}\in\Z$。注意 $\nu(-1)$ 通过 $K(\mathcal{M}_{coh}(\D_X,N_0))$ 分解，且由 2.4 有
$$\nu(I_{O_D})(-1)=\delta_D(-1)\qquad(D\in W_\Theta\backslash W).$$
于是
$$\nu(L_{O_C})(-1)=\sum_{D\le C}\lambda_{CD}\nu(I_{O_D})(-1)=\sum_{D\le C}\lambda_{CD}\delta_D(-1).$$
故由 $P_{CD}$ 的定义可得 $\lambda_{CC}=1$ 且 $P_{CD}(-1)=\lambda_{CD}$。
\end{proof}

\begin{annote}
底本此处把 $\nu(I_D)(-1)=\delta_D(-1)$ 的依据引作 2.6。但 2.6 是关于 $\chi(\calO(\lambda+\rho))$ 的推论，与该式无关；给出该式的是 2.4 之后的计算（用 1.8 得 $\nu(I(O,\lambda))=\sum_{v\in W_\Theta}q^{\ell(v)}\delta_{vw}=\delta_C$）。已按 2.4 读出。[来源：底本 p.119--121，2.4 与 2.6 的陈述]
\end{annote}

这就给出一个计算无穷小特征标 $\chi_\rho$ 下广义 Verma 模中不可约模重数的有效算法。可取 $W_\Theta\backslash W$ 上一个与 Bruhat 序相容的序关系。则矩阵 $(\lambda_{CD};C,D\in W_\Theta\backslash W)$ 是下三角的，且对角元为 $1$。设 $(\mu_{CD};C,D\in W_\Theta\backslash W)$ 为其逆矩阵的系数。由 3.6 可见
{\small
\begin{align*}
\chi(I_{O_C})&=\sum_{E\in W_\Theta\backslash W}\sum_{D\in W_\Theta\backslash W}\mu_{CD}\lambda_{DE}\chi(I_{O_E})
=\sum_{D\in W_\Theta\backslash W}\mu_{CD}\Bigl(\sum_{E\in W_\Theta\backslash W}\lambda_{DE}\chi(I_{O_E})\Bigr)\\
&=\sum_{D\in W_\Theta\backslash W}\mu_{CD}\chi(L_{O_D})
=\sum_{D\le C}\mu_{CD}\chi(L_{O_D}),
\end{align*}
}
且对任意 $C\in W_\Theta\backslash W$ 有 $\mu_{CC}=1$。于是由 2.12、2.5 与 V.1.19，最终得到下述结果。

\begin{corollary}\label{cor:06.3.7}
不可约模 $L(-v\rho)$（$v\in{}^{\Theta}W$）在广义 Verma 模 $M_\Theta(-w\rho)$（$w\in{}^{\Theta}W$）中的重数，等于 $\mu_{CD}$，其中 $C,D\in W_\Theta\backslash W$ 分别是 $w$ 与 $v$ 所在的陪集。
\end{corollary}
